\begin{frame}{skip-gram의 실제 구조}
  ``작은 신경망''을 정확히 보면:
  \small
  \begin{tabular}{@{}lp{0.62\textwidth}@{}}
    \toprule
    층 & 크기 \\
    \midrule
    입력층 & 어휘 크기 $V$개 (\kb{one-hot} 입력) \\
    은닉층 & $d$개 (= \textbf{임베딩 차원}) \\
    출력층 & 다시 어휘 크기 $V$개 \\
    \bottomrule
  \end{tabular}
  \vskip1em
  \begin{itemize}
    \item 중심 단어 $c$의 임베딩 = $\boldsymbol{v}_{c}$ (은닉층의
      ``가중치 행''). \; 출력층에서 단어 $t$를 가리키는 벡터 =
      $\boldsymbol{u}_{t}$.
    \item ``중심 $c$의 바로 옆에 $t$가 나올 확률'' (2.3절 시그모이드):
  \end{itemize}
  \[
    p(t \mid c) = \sigma(\boldsymbol{u}_{t}^{\top}\boldsymbol{v}_{c})
    = \frac{1}{1 + e^{-\boldsymbol{u}_{t}^{\top}\boldsymbol{v}_{c}}}
  \]
  즉, \textbf{두 벡터의 내적이 클수록} $t$가 $c$ 옆에 나올
  가능성이 크다고 학습됨.
\end{frame}
